\documentclass[aspectratio=169,10pt,english]{beamer}

\input{../../../_preambulo_beamer.tex}
\input{../../../_comandos_beamer.tex}

\selectlanguage{english}

\title{MA1001B - Statistical Modeling for Decision Making}
\subtitle{Lesson 31: Chi-Square and a Confidence Interval for Variance}
\author{}
\institute{Tecnol\'ogico de Monterrey}
\date{}

\begin{document}

\begin{frame}[plain]
  \vspace{1.2cm}
  {\Large\bfseries MA1001B - Statistical Modeling for Decision Making\par}
  \vspace{0.7cm}
  {\large Lesson 31: Chi-Square and a Confidence Interval for Variance\par}
  \vspace{0.7cm}
  {\color{TecRojo}\rule{\textwidth}{1.2pt}}
  \vspace{0.8cm}

  MA1001B Faculty\\[0.3cm]
  Term\\[0.3cm]
  Tecnol\'ogico de Monterrey
\end{frame}

\begin{frame}{The Variance-Interval Question}
  \begin{alertblock}{Guiding question}
    How is a 95\% interval for an unknown $\sigma^{2}$ built from $s^{2}$
    when the population is treated as normal?
  \end{alertblock}

  \vspace{0.6em}
  By the end of this lesson, you can:
  \begin{itemize}
    \item compute $s^{2}$ with divisor $n-1$;
    \item obtain chi-square critical values with $\mathrm{df}=n-1$;
    \item form $((n-1)s^{2}/\chi^{2}_{U},(n-1)s^{2}/\chi^{2}_{L})$;
    \item explain why non-normal data invalidate the interval.
  \end{itemize}

  \vfill
  \scriptsize All cycle times used today are synthetic. No real company data are used.
\end{frame}

\begin{frame}{A Small Normal Sample of Cycle Times}
  \small
  \begin{center}
    \begin{tabular}{lr}
      \toprule
      \textbf{Quantity} & \textbf{Value} \\
      \midrule
      Sample size $n$ & 20 cycles \\
      Degrees of freedom $n-1$ & 19 \\
      Sample $s^{2}$ (seed 42) & 6.814885 \\
      Hidden $\sigma^{2}$ & 9 \\
      \bottomrule
    \end{tabular}
  \end{center}

  \vspace{0.5em}
  \begin{exampleblock}{Pivot}
    If the population is normal, $(n-1)s^{2}/\sigma^{2}\sim\chi^{2}_{n-1}$.
    That pivot is inverted to obtain an interval for $\sigma^{2}$.
  \end{exampleblock}
\end{frame}

\begin{frame}[fragile]{Step 1: Sample Variance}
  \begin{columns}[T]
    \begin{column}{0.42\textwidth}
      \small
      \begin{lstlisting}
s2 = np.var(sample, ddof=1)
df = n - 1
      \end{lstlisting}

      \begin{block}{Verified output}
        $\mathrm{df}=19$\\
        $s=2.610533$\\
        $s^{2}=6.814885$
      \end{block}
    \end{column}
    \begin{column}{0.58\textwidth}
      \centering
      \includegraphics[width=\textwidth]{../figures/chi2_var_01_sample_variance.png}
      \vspace{0.2em}
      \scriptsize Cycle-time histogram from Step 1
    \end{column}
  \end{columns}
\end{frame}

\begin{frame}{Invert the Chi-Square Pivot}
  \small
  Critical values from \texttt{scipy.stats.chi2.ppf}:
  \[
    \chi^{2}_{0.025,19}=8.906516,\qquad
    \chi^{2}_{0.975,19}=32.852327.
  \]
  Interval:
  \[
    \left[\frac{19\times 6.814885}{32.852327},
          \frac{19\times 6.814885}{8.906516}\right]
    =[3.941359, 14.537986].
  \]
  The larger chi-square value produces the lower bound. This seed-42 interval
  covers the hidden $\sigma^{2}=9$.
\end{frame}

\begin{frame}[fragile]{Step 2: The 95\% Chi-Square Interval}
  \begin{columns}[T]
    \begin{column}{0.40\textwidth}
      \small
      \begin{lstlisting}
lower = df * s2 / chi2_u
upper = df * s2 / chi2_l
      \end{lstlisting}

      \begin{block}{Verified output}
        $\chi^{2}_{L}=8.906516$\\
        $\chi^{2}_{U}=32.852327$\\
        CI $[3.941359, 14.537986]$
      \end{block}
    \end{column}
    \begin{column}{0.60\textwidth}
      \centering
      \includegraphics[width=\textwidth]{../figures/chi2_var_02_variance_interval.png}
      \vspace{0.2em}
      \scriptsize Variance interval from Step 2
    \end{column}
  \end{columns}
\end{frame}

\begin{frame}{Step 3: Non-Normal Data Invalidate the Interval}
  \begin{columns}[T]
    \begin{column}{0.42\textwidth}
      \small
      5000 replications, $n=20$.

      \vspace{0.4em}
      Normal cycles: coverage $0.949400$.

      \vspace{0.4em}
      Exponential cycles: coverage $0.737200$.

      \vspace{0.5em}
      \textbf{Limit:} the chi-square interval for $\sigma^{2}$ assumes a
      normal population. Skewed cycle times destroy the nominal 95\%.
    \end{column}
    \begin{column}{0.58\textwidth}
      \centering
      \includegraphics[width=\textwidth]{../figures/chi2_var_03_nonnormal_limit.png}
      \vspace{0.2em}
      \scriptsize Empirical coverage from Step 3
    \end{column}
  \end{columns}
\end{frame}

\begin{frame}{Concept Check}
  \begin{enumerate}
    \item Why does $s^{2}$ use divisor $n-1$ rather than $n$?
    \item Which chi-square quantile sits in the denominator of the lower
      bound?
    \item Is the interval symmetric about $s^{2}=6.814885$?
    \item What happens to coverage when cycle times are exponential?
  \end{enumerate}

  \vspace{0.6em}
  \begin{block}{Connection to Lesson 32}
    The session can continue directly into the $F$ distribution and a
    comparison of two variances.
  \end{block}

  \vfill
  \scriptsize Source: OpenStax, \textit{Introductory Business Statistics 2e},
  Chapter 11, Section 11.2. CC BY-NC-SA.
\end{frame}

\appendix



\end{document}
